\textbf{Cómo localizar los componentes del expediente.}\par
\begin{description}
\item[Fuentes] \ruta{main.tex} contiene la estructura institucional. Los bloques técnicos se incluyen desde \ruta{resultados/}. La bibliografía utiliza \ruta{sample.bib}.
\item[Diagramas y código] \ruta{diagramas/} conserva las fuentes TikZ editables. \ruta{codigo/} contiene los fragmentos citados con numeración respecto al commit auditado.
\item[Evidencia visual] \ruta{capturas/} conserva las capturas de ejecución y su manifiesto en \ruta{evidencias/capturas.md}. Los wireframes, fuentes de referencia de diseño, se guardan en \ruta{wireframes/}.
\item[Evidencia automatizada] \ruta{evidencias/auditoria-final.trx} contiene la ejecución de auditoría; se incluyen también las comparaciones de regresión y el registro de SQLite nueva.
\item[Estado de entrega] Se conservaron siete capturas auténticas con cursor o halo, sin alterar datos ni resultados funcionales. El registro de cierre visual documenta su procedencia y el alcance de las comprobaciones DPI.
\end{description}
